\chapter{Lý Do Bá Đạo Và Cách Đáp}
\chapterintro{Gặp lý do lạ không cần mặt nặng ngay}{Giữ được sự bình tĩnh và một chút duyên là giữ được thế chủ động.}

\section{Em không làm bài vì hôm qua em buồn}
\begin{manthoai}{Màn thoại}
\textbf{Học sinh:} Em không làm bài vì hôm qua em buồn.\\
\textbf{Thầy/Cô:} Thầy cô hiểu. Giờ mình làm phần ngắn nhất trước để đỡ buồn thêm nhé.
\end{manthoai}
\begin{dungkhinào}
Dùng khi muốn thừa nhận cảm xúc của học sinh nhưng không bỏ luôn trách nhiệm học tập.
\end{dungkhinào}
\begin{nhipthem}
Đáp bằng việc cụ thể để kéo học sinh ra khỏi lý do vòng vo.
\end{nhipthem}
\ngancach

\section{Em tưởng hôm nay không kiểm tra bài cũ}
\begin{manthoai}{Màn thoại}
\textbf{Học sinh:} Em tưởng hôm nay không kiểm tra bài cũ.\\
\textbf{Thầy/Cô:} Tưởng tượng của em phong phú đấy. Giờ mình quay lại phần thực tế nhé.
\end{manthoai}
\begin{dungkhinào}
Hợp khi lý do mang tính chống đỡ rõ nhưng chưa cần làm căng.
\end{dungkhinào}
\begin{nhipthem}
Sau câu đáp vui, xử lý việc chính thật gọn để tiết học không bị trôi vì đối đáp quá lâu.
\end{nhipthem}
\ngancach

\section{Em làm rồi nhưng quên mang}
\begin{manthoai}{Màn thoại}
\textbf{Học sinh:} Em làm rồi nhưng quên mang.\\
\textbf{Thầy/Cô:} Vậy hôm nay em mang trí nhớ ra chứng minh trước nhé.
\end{manthoai}
\begin{dungkhinào}
Hợp khi muốn chuyển từ lời biện hộ sang một cách kiểm tra thay thế ngay tại lớp.
\end{dungkhinào}
\begin{nhipthem}
Mời em trả lời ngắn, nêu ý chính hoặc làm một phần nhỏ thay cho vở.
\end{nhipthem}
\ngancach

\section{Em quên bài vì tối qua em ngủ quên ạ}
\begin{manthoai}{Màn thoại}
\textbf{Học sinh:} Em quên bài vì tối qua em ngủ quên ạ.\\
\textbf{Thầy/Cô:} Thầy cô tôn trọng giấc ngủ của em. Giờ mình thức dậy bằng một ý chính trước nhé.
\end{manthoai}
\begin{dungkhinào}
Hợp với những lý do nửa thật nửa đỡ đòn, khi thầy cô muốn kéo nhanh về việc học thay vì tranh cãi.
\end{dungkhinào}
\begin{nhipthem}
Sau câu đáp, cho em làm một phần thật ngắn để chuyển từ lý do sang hành động.
\end{nhipthem}
\ngancach

\section{Em đang nộp lý do hay nộp giải pháp?}
\begin{manthoai}{Màn thoại}
\textbf{Thầy/Cô:} Em đang nộp lý do hay nộp giải pháp?\\
\textbf{Học sinh:} Dạ giải pháp ạ.\\
\textbf{Thầy/Cô:} Tốt. Vậy nói thầy cô nghe mình xử lý thiếu sót này thế nào.
\end{manthoai}
\begin{dungkhinào}
Dùng khi học sinh đang kéo dài phần biện minh và lớp bắt đầu mất nhịp vì màn giải thích.
\end{dungkhinào}
\begin{nhipthem}
Câu này hiệu quả vì nó đẩy đối thoại từ chống đỡ sang chịu trách nhiệm.
\end{nhipthem}
